\section{Median of a Row-wise Sorted Matrix}
\label{sec:bs-matrix-median}

\problemstatement{We are given an $n \times m$ matrix in which every row
is sorted in increasing order (with no chaining guarantee between rows),
and $nm$ is odd. Find the median of all $nm$ elements.}

\begin{patternbox}
This closing problem of the chapter combines two ideas from earlier
sections: \emph{binary search on the answer}
(Sections~\ref{sec:bs-sqrt}--\ref{sec:bs-kth-element}), where the
candidate space is a range of \emph{values} rather than indices, and
\emph{per-row binary search} (Section~\ref{sec:bs-row-max-ones}), where
row-sortedness is exploited independently per row inside a helper
function. The keyword combination \emph{``median''} plus \emph{``each row
sorted''} plus a request for better than $O(nm \log(nm))$ (which a full
sort-and-pick would cost) should trigger exactly this combination.
\end{patternbox}

\subsection*{Understanding the problem carefully}

The median of $nm$ elements (with $nm$ odd) is the value $v$ such that
exactly $(nm-1)/2$ elements are strictly smaller than $v$, and $v$ itself
is present in the matrix. Equivalently, $v$ is the smallest value for
which $\textit{countLessOrEqual}(v) \ge \lceil nm/2 \rceil$, where
$\textit{countLessOrEqual}(v)$ counts how many matrix elements are $\le
v$. We do not need to merge or sort the matrix to compute this count: for
each row, since the row is independently sorted, the count of elements
$\le v$ in that row is exactly $\textit{upperBound}(\text{row}, v)$
(Section~\ref{sec:bs-upper-bound}, the first index strictly greater than
$v$, which equals the count of elements $\le v$). Summing this across all
$n$ rows gives $\textit{countLessOrEqual}(v)$ for the whole matrix in
$O(n \log m)$.

\begin{ideabox}[Candidate Space and Predicate]
The candidate space is integer values in
$[\min(\textit{matrix}), \max(\textit{matrix})]$ (the minimum and maximum
single elements across the whole matrix, found by checking the first and
last elements of each row, since rows are sorted). The predicate is
$P(v) = (\textit{countLessOrEqual}(v) \ge \lceil nm/2 \rceil)$, and we want
the smallest $v$ where $P$ holds.
\end{ideabox}

\subsection*{Why this nested binary search correctly finds the median}

As $v$ increases, $\textit{countLessOrEqual}(v)$ can only increase or stay
the same (a larger threshold never excludes an element that a smaller
threshold included), so $P$ is monotonic, justifying the outer binary
search over $v$. The smallest $v$ satisfying $P$ is, by definition, the
smallest value such that at least half of the matrix's elements do not
exceed it -- which is exactly the definition of the median for an odd
total count. Crucially, $v$ itself does not need to be searched for
separately within the matrix: because $\textit{countLessOrEqual}$ jumps by
at least $1$ at $v$ if and only if $v$ is actually present in the matrix
(if $v$ were absent, the count at $v$ and at $v-1$ would be identical, so
the crossover could not happen exactly at an absent value), the binary
search automatically lands on a value that is genuinely in the matrix.

\subsection*{Algorithm}

\begin{enumerate}
  \item $\textit{lo} \gets \min$ over all rows' first elements;
        $\textit{hi} \gets \max$ over all rows' last elements.
  \item While $\textit{lo} \le \textit{hi}$:
  \begin{enumerate}
    \item $\textit{mid} \gets \textit{lo}+(\textit{hi}-\textit{lo})/2$.
    \item For each row, add $\textit{upperBound}(\text{row},
          \textit{mid})$ to get $\textit{count}$.
    \item If $\textit{count} < \lceil nm/2 \rceil$: $\textit{lo} \gets
          \textit{mid}+1$.
    \item Else: $\textit{hi} \gets \textit{mid}-1$.
  \end{enumerate}
  \item Return \textit{lo}.
\end{enumerate}

\begin{algorithm}[H]
\caption{Median of a Row-wise Sorted Matrix}
\begin{algorithmic}[1]
\Procedure{MatrixMedian}{\textit{matrix}}
  \State $\textit{lo} \gets \min_i \textit{matrix}[i][0]$, $\textit{hi} \gets \max_i \textit{matrix}[i][m-1]$
  \State $\textit{half} \gets \lceil nm/2 \rceil$
  \While{$\textit{lo} \le \textit{hi}$}
    \State $\textit{mid} \gets \textit{lo}+(\textit{hi}-\textit{lo})/2$
    \State $\textit{count} \gets \sum_i \Call{UpperBound}{\textit{matrix}[i], \textit{mid}}$
    \If{$\textit{count} < \textit{half}$}
      \State $\textit{lo} \gets \textit{mid} + 1$
    \Else
      \State $\textit{hi} \gets \textit{mid} - 1$
    \EndIf
  \EndWhile
  \State \Return \textit{lo}
\EndProcedure
\end{algorithmic}
\end{algorithm}

\subsection*{Dry run}

\dryrun{Take $\textit{matrix} = \begin{bmatrix} 1&3&5 \\ 2&6&9 \\
3&6&9 \end{bmatrix}$ ($n=m=3$, $nm=9$).}

$\textit{half} = \lceil 9/2 \rceil = 5$. $\textit{lo}=\min(1,2,3)=1$,
$\textit{hi}=\max(5,9,9)=9$.

\begin{itemize}
  \item $\textit{lo}=1,\textit{hi}=9$: $\textit{mid}=5$. Counts $\le 5$:
        row $1$: $[1,3,5]$, upper bound of $5$ is index $3$ (all $3$
        qualify). Row $2$: $[2,6,9]$, upper bound of $5$ is index $1$
        (just the $2$). Row $3$: $[3,6,9]$, upper bound of $5$ is index
        $1$ (just the $3$). Total $=3+1+1=5 \ge 5$. $\textit{hi}=4$.
  \item $\textit{lo}=1,\textit{hi}=4$: $\textit{mid}=2$. Counts $\le 2$:
        row $1$: upper bound of $2$ is index $1$ (just the $1$). Row $2$:
        upper bound of $2$ is index $1$ (just the $2$). Row $3$: upper
        bound of $2$ is index $0$ (none). Total $=1+1+0=2 < 5$.
        $\textit{lo}=3$.
  \item $\textit{lo}=3,\textit{hi}=4$: $\textit{mid}=3$. Counts $\le 3$:
        row $1$: index $2$ (the $1,3$). Row $2$: index $1$ (just the
        $2$). Row $3$: index $1$ (just the $3$). Total $=2+1+1=4 < 5$.
        $\textit{lo}=4$.
  \item $\textit{lo}=4,\textit{hi}=4$: $\textit{mid}=4$. Counts $\le 4$:
        same as $\le 3$ for these rows (no element equals $4$), so total
        $=4 < 5$. $\textit{lo}=5$.
  \item $\textit{lo}=5 > \textit{hi}=4$: loop ends.
\end{itemize}

The answer is $\textit{lo} = 5$. Checking by hand: sorting all $9$
elements gives $[1,2,3,3,5,6,6,9,9]$, and the middle ($5$th) element is
indeed $5$.

\subsection*{C++ implementation}

\begin{lstlisting}
#include <bits/stdc++.h>
using namespace std;

int upperBoundRow(vector<int>& row, int X) {
    int lo = 0, hi = row.size() - 1, ans = row.size();
    while (lo <= hi) {
        int mid = lo + (hi - lo) / 2;
        if (row[mid] > X) { ans = mid; hi = mid - 1; }
        else lo = mid + 1;
    }
    return ans;
}

int matrixMedian(vector<vector<int>>& matrix) {
    int n = matrix.size(), m = matrix[0].size();
    int lo = INT_MAX, hi = INT_MIN;
    for (auto& row : matrix) {
        lo = min(lo, row[0]);
        hi = max(hi, row[m - 1]);
    }

    int half = (n * m) / 2 + 1; // ceil(nm/2) for odd nm

    while (lo <= hi) {
        int mid = lo + (hi - lo) / 2;
        int count = 0;
        for (auto& row : matrix) {
            count += upperBoundRow(row, mid);
        }

        if (count < half) {
            lo = mid + 1;
        } else {
            hi = mid - 1;
        }
    }
    return lo;
}

int main() {
    vector<vector<int>> matrix = {
        {1, 3, 5},
        {2, 6, 9},
        {3, 6, 9}
    };
    cout << "Median: " << matrixMedian(matrix) << endl;
    return 0;
}
\end{lstlisting}

\begin{complexitybox}
\textbf{Time Complexity.} The candidate value range halves each iteration
($O(\log(\max - \min))$ iterations), and each feasibility check sums a
per-row binary search costing $O(\log m)$ across $n$ rows, giving an
overall time complexity of
\[
  O(n \log m \log(\max-\min)),
\]
substantially better than sorting all $nm$ elements directly.
\textbf{Space Complexity.} $O(1)$ auxiliary space.
\end{complexitybox}

\begin{takeawaybox}
This closing problem is a deliberate capstone: it requires recognizing
that ``binary search on the answer'' (search over a value range using a
monotonic count predicate) and ``per-row binary search'' (exploiting
row-sortedness inside the predicate) are not mutually exclusive
techniques -- they nest naturally, with the outer search handling
\emph{which} value to test and the inner per-row searches handling
\emph{how cheaply} to test it. Recognizing when two patterns from
different parts of this chapter can be combined, rather than always
reaching for a single isolated template, is the most advanced skill this
chapter has to offer.
\end{takeawaybox}
